\subsection{Newton閾値と境界matrix factorizationの非同値性}
\label{subsec:boundary-thresholds-mf}
前の定理で得た二枝のNewton面を、乗数イデアルと局所特異点圏に適用する。
$n\ge1$、$b=b_n,c=c_n$、$\Lambda=3+2\sqrt2$と略記し、
式\eqref{eq:boundary-conic-condition}を仮定する。
定理\ref{thm:two-boundary-Puiseux}より、唯一の境界特異点は$p_w$であり、
局所被約方程式$f_n$は重み$(b,c)$の最低部分
\[
 (y^b-\rho_+x^c)(y^b-\rho_-x^c)
 =y^{2b}-2\kappa x^cy^b+x^{2c}
\]
を持つ。$\rho_+\ne\rho_-$、$\rho_+\rho_-=1$なのでNewton非退化である。
ここでの不変量は\emph{閉包曲線の境界特異点}のものであり、
曲面$U$そのものや元の非可換H3圏の不変量ではない。

\begin{theorem}[局所対数標準閾値と乗数イデアル]
\label{thm:boundary-lct}
$0<\tau<1$について、$p_w$での乗数イデアルは
\[
 \mathcal J(S,\tau\overline C_n)_{p_w}
 =\left\langle x^iy^j:
 b(i+1)+c(j+1)>2bc\tau\right\rangle
 \subset\mathcal O_{S,p_w}^{\rm an}.
\]
特に
\begin{equation}\label{eq:pell-lct}
 \boxed{\operatorname{lct}_{p_w}(S;\overline C_n)
 =\frac{b+c}{2bc}},\qquad
 \lim_{n\to\infty}-\frac1n\log\operatorname{lct}_{p_w}(S;\overline C_n)
 =\log(3+2\sqrt2).
\end{equation}
ここでは境界$D$に係数1を与えておらず、対$(S,D)$に対する相対閾値ではない。
\end{theorem}
\begin{proof}
Newton非退化性により、前の二枝のNewton面に適合するトーリック解消は$f_n$のlog resolutionでもある。
正の原始ray$(r,s)$に対応する例外曲線では、log discrepancyは$r+s$、
$f_n$の位数は$\min(2cr,2bs)$である。
$\tau<1$ではstrict transformは乗数イデアルへ条件を加えず、
各単項式の条件は全rayについて
$(i+1)r+(j+1)s>\tau\min(2cr,2bs)$となる。
Newton多面体の唯一のコンパクト面に対応する$(r,s)=(b,c)$の条件が
他の条件を含意し、表示式を得る。
これは非退化主部に対するHowaldのトーリック計算法の局所版を直接適用したものでもある。\footnote{
J. Howald, \emph{Multiplier Ideals of Sufficiently General Polynomials},
\href{https://arxiv.org/abs/math/0303203}{arXiv:math/0303203}, Theorems 11--12 and Corollary 13.
本稿の局所冪級数では、上記のNewton面と明示的トーリック解消を使って計算している。}
定数$1$が初めて乗数イデアルから外れる値が$(b+c)/(2bc)$である。
$b_n,c_n$はともに正定数倍の$(3+2\sqrt2)^n$に漸近するので極限式も従う。
\end{proof}

\begin{corollary}[境界matrix factorizationのHochschild次元]
\label{cor:boundary-mf-growth}
$\mathcal M_n=\operatorname{MF}(\mathbb C[[x,y]],f_n)$を
上で選んだ局所被約方程式の$\mathbb Z/2$-graded dg圏とする。
このとき
\[
 \mathrm{HH}^{\bar0}(\mathcal M_n)
 \cong\frac{\mathbb C[[x,y]]}{(\partial_x f_n,\partial_y f_n)},
 \qquad\mathrm{HH}^{\bar1}(\mathcal M_n)=0,
\]
\[
 \dim\mathrm{HH}^{\bar0}(\mathcal M_n)
 =(2b_n-1)(2c_n-1),\qquad
 \lim_{n\to\infty}\frac1n\log\dim\mathrm{HH}^{\bar0}(\mathcal M_n)
 =2\log(3+2\sqrt2).
\]
特に異なる$n\ge1$の$\mathcal M_n$は$\mathbb Z/2$-graded dg Morita同値ではない。
\end{corollary}
\begin{proof}
$f_n$は孤立した臨界点を持ち、形式二変数環はregularである。
従ってDyckerhoffの既知の計算\footnote{
T. Dyckerhoff, \emph{Compact generators in categories of matrix factorizations},
\href{https://arxiv.org/abs/0904.4713v5}{arXiv:0904.4713v5},
Corollary 6.5（有限rankの二周期dg圏）。
Theorem 6.6はhomology側の対応であり、ここではcohomology側を使用する。}
を適用できる。そのJacobian環の次元は局所Milnor数なので、定理\ref{thm:boundary-torus-link}から上の公式を得る。
$b_n,c_n$はともに真に増大し、Hochschild cohomologyはdg Morita不変だから
異なる$n$の同値は存在しない。
\end{proof}

ここで$\mathcal M_n$は\emph{閉包曲線の境界特異点}に付随する圏である。
これを$\operatorname{Perf}(U)$、元の非可換H3圏、あるいは
$\operatorname{Perf}(S)\to\operatorname{Perf}(U)$の核と同一視したわけではない。
同じ曲線のアフィン部分は全て$\mathbb G_m$である一方、
境界の特異点圏はHochschild次元だけでも相異なる。
この圏の\emph{列}の次元成長も、単一の自己関手のentropyではない。

\begin{remark}[同じアフィン曲線と異なる境界の圏]
最初の三例では
\[
\begin{array}{c|c|c|c}
 n&(b_n,c_n)&\dim\mathrm{HH}^{\bar0}(\mathcal M_n)&\operatorname{lct}\\ \hline
 1&(2,5)&27&7/20\\
 2&(12,29)&1311&41/696\\
 3&(70,169)&46843&239/23660
\end{array}
\]
となる。アフィン部分$C_n$は全て$\mathbb G_m$で、
通常のアフィン対象複雑度は有界であるにもかかわらず、
境界のMF圏はHochschild次元だけで互いに区別できる。
一方、局所Milnor monodromyは各$n$で有限位数である。
これらの異なる意味の成長を単一の自己関手のentropyにまとめない。

HowaldとDyckerhoffの一般理論は既知である。
本稿では指定したPell反復族に適用し、閾値の減衰率$\log\Lambda$と
局所圏の次元増大率$2\log\Lambda$を得た。
\texttt{verify\_boundary\_thresholds.py}はLaurent帰納指数、整数行列、
半群・固有値多重集合と有限範囲の数値を厳密に再現する。
全$n$の理論とMF定理を有限実験だけで証明したとはせず、
個別の適用の新規性も網羅的には確認していない。
\end{remark}
